\documentclass[10pt,a4paper]{amsart}
\usepackage{amsmath,amssymb,amsthm,graphicx}
% Retained minimal preamble for the Stacks Project chapter more-algebra.tex (More on Algebra).
% The chapter file pinned-original-source.tex carries only \input{preamble} on its first line, so no
% preamble text lives inside the chapter itself.  The Stacks Project uses one single shared preamble for
% every chapter; the snapshot stages it identically as sources/stacks-chow/preamble.tex and
% sources/stacks-etale-cohomology/preamble.tex.  The block below is copied mechanically, line for line,
% from sources/stacks-chow/preamble.tex lines 14-18 (the Xy-pic setup) and 154-217 (the native theorem
% declarations together with the mathematical macros).  Nothing was added, removed or redefined; the
% wrapper class and the amsmath/amssymb/amsthm/graphicx base packages are supplied by the assembler.
% Chapter: more-algebra.tex -- More on Algebra.
% For commutative diagrams we use Xy-pic
\usepackage[all]{xy}

% We use 2cell for 2-commutative diagrams.
\xyoption{2cell}

% Theorem environments.
%
\theoremstyle{plain}
\newtheorem{theorem}[subsection]{Theorem}
\newtheorem{proposition}[subsection]{Proposition}
\newtheorem{lemma}[subsection]{Lemma}

\theoremstyle{definition}
\newtheorem{definition}[subsection]{Definition}
\newtheorem{example}[subsection]{Example}
\newtheorem{exercise}[subsection]{Exercise}
\newtheorem{situation}[subsection]{Situation}

\theoremstyle{remark}
\newtheorem{remark}[subsection]{Remark}
\newtheorem{remarks}[subsection]{Remarks}

\numberwithin{equation}{subsection}

% Macros
%
\def\lim{\mathop{\mathrm{lim}}\nolimits}
\def\colim{\mathop{\mathrm{colim}}\nolimits}
\def\Spec{\mathop{\mathrm{Spec}}}
\def\Hom{\mathop{\mathrm{Hom}}\nolimits}
\def\Ext{\mathop{\mathrm{Ext}}\nolimits}
\def\SheafHom{\mathop{\mathcal{H}\!\mathit{om}}\nolimits}
\def\SheafExt{\mathop{\mathcal{E}\!\mathit{xt}}\nolimits}
\def\Sch{\mathit{Sch}}
\def\Mor{\mathop{\mathrm{Mor}}\nolimits}
\def\Ob{\mathop{\mathrm{Ob}}\nolimits}
\def\Sh{\mathop{\mathit{Sh}}\nolimits}
\def\NL{\mathop{N\!L}\nolimits}
\def\CH{\mathop{\mathrm{CH}}\nolimits}
\def\proetale{{pro\text{-}\acute{e}tale}}
\def\etale{{\acute{e}tale}}
\def\QCoh{\mathit{QCoh}}
\def\Ker{\mathop{\mathrm{Ker}}}
\def\Im{\mathop{\mathrm{Im}}}
\def\Coker{\mathop{\mathrm{Coker}}}
\def\Coim{\mathop{\mathrm{Coim}}}

% Boxtimes
%
\DeclareMathSymbol{\boxtimes}{\mathbin}{AMSa}{"02}

%
% Macros for moduli stacks/spaces
%
\def\QCohstack{\mathcal{QC}\!\mathit{oh}}
\def\Cohstack{\mathcal{C}\!\mathit{oh}}
\def\Spacesstack{\mathcal{S}\!\mathit{paces}}
\def\Quotfunctor{\mathrm{Quot}}
\def\Hilbfunctor{\mathrm{Hilb}}
\def\Curvesstack{\mathcal{C}\!\mathit{urves}}
\def\Polarizedstack{\mathcal{P}\!\mathit{olarized}}
\def\Complexesstack{\mathcal{C}\!\mathit{omplexes}}
% \Pic is the operator that assigns to X its picard group, usage \Pic(X)
% \Picardstack_{X/B} denotes the Picard stack of X over B
% \Picardfunctor_{X/B} denotes the Picard functor of X over B
\def\Pic{\mathop{\mathrm{Pic}}\nolimits}
\def\Picardstack{\mathcal{P}\!\mathit{ic}}
\def\Picardfunctor{\mathrm{Pic}}
\def\Deformationcategory{\mathcal{D}\!\mathit{ef}}

% --- appended 2026-10-12 (Chief): the project own definitions of the reference/slogan/history
% display environments, copied mechanically from the shared Stacks preamble lines 8-12
% (\usepackage{verbatim} plus the three \newenvironment lines, which the chapter preamble keeps
% verbatim from the project preamble). Nothing added, removed or redefined: these environments
% carry no mathematical content and the project itself typesets them as comments. Without these
% three lines a counted statement whose own text uses them fails with Environment slogan undefined.
\usepackage{verbatim}
\newenvironment{reference}{\comment}{\endcomment}
\newenvironment{slogan}{\comment}{\endcomment}
\newenvironment{history}{\comment}{\endcomment}
\usepackage{cleveref}
\crefname{originallemma}{original lemma}{original lemmas}
\Crefname{originallemma}{Original lemma}{Original lemmas}
\makeatletter
\newlabel{algebra-lemma-integral-sub-dim-equal}{{00OK}{source}{}{sourcepointer-1}{}}
\newlabel{algebra-lemma-integral-sub-dim-equal@cref}{{[originallemma][0][]00OK}{[1][0][]source}{}{}{}}
\makeatother
\crefname{originallemma}{original lemma}{original lemmas}
\Crefname{originallemma}{Original lemma}{Original lemmas}
\makeatletter
\newlabel{algebra-lemma-krull-akizuki}{{00PG}{source}{}{sourcepointer-2}{}}
\newlabel{algebra-lemma-krull-akizuki@cref}{{[originallemma][0][]00PG}{[1][0][]source}{}{}{}}
\makeatother
\crefname{originallemma}{original lemma}{original lemmas}
\Crefname{originallemma}{Original lemma}{Original lemmas}
\makeatletter
\newlabel{algebra-lemma-finite-extension-residue-fields-dimension-1}{{031F}{source}{}{sourcepointer-3}{}}
\newlabel{algebra-lemma-finite-extension-residue-fields-dimension-1@cref}{{[originallemma][0][]031F}{[1][0][]source}{}{}{}}
\makeatother
\title{Noetherianity, surjectivity and finite residue extensions for the integral closure in a base field extension}
\author{The Stacks Project Authors}
\date{stacks-project:more-algebra.tex:local-snapshot-20261012; GNU FDL 1.2 or later}
\begin{document}
\maketitle
\setcounter{subsection}{1}
\noindent\textit{Editorial note.} All delivered mathematical text is the author's own, copied from the pinned Stacks Project chapter file more-algebra.tex; no statement or proof step was rewritten, reordered or removed. The counted statement is the labeled lemma environment at source lines 35925-35935 (canonical tag 09IH) and the counted proof is the proof environment at source lines 35937-35958, which immediately follows it and is its own proof.\par\medskip
\noindent\textit{Editorial note.} The printed text writes suboject in the section introduction rather than in this proof; inside the counted spans the wording is copied exactly as printed and no correction is applied.\par\medskip
\noindent\textit{Editorial note.} The preamble span runs to line 90, the current end of the retained preamble, which includes the appended environment definitions; the counted text does not use them, and the span is taken whole so that the file version is unambiguous.\par\medskip
\noindent\textit{Editorial note.} The source URL is the source-level Stacks Project URL, as required when more than one card is produced from one source; the per-claim tag is carried by claim\_key and by this note. The counted statement and proof cite no bibliography entry, so this card needs no bibliography.\par\medskip
\noindent\textit{Allowed standard prerequisites (editorial).}\begin{itemize}
\item A localization of a finite type algebra at a prime is integral over the localization of the finite type algebra after an integral base change, so dimensions agree. used at proof lines 35937-35945. Shows that the localized finite type algebra over the new base is one dimensional and Noetherian, which is the setting for the integral closure.
\item Krull-Akizuki: the integral closure of a one dimensional Noetherian domain in a finite extension of its fraction field is Noetherian. invoked at proof lines 35948-35952. Makes each of the finitely many integral closures of the quotients Noetherian.
\item For a one dimensional local ring the residue field extension in a finite extension of the fraction field is finite. invoked at proof lines 35952-35955. Gives the finiteness of the residue field extensions asserted in the statement.
\end{itemize}\medskip
\ifdefined\hypertarget\hypertarget{sourcepointer-1}{}\fi
\noindent\textit{Source reference.} Original lemma 00OK: sources/stacks-chow/tags.txt, source lines 898--898. The reference is applied at proof lines 35944-35945 to conclude that the localized ring has dimension one, dimensions being preserved under integral extensions; the list of minimal primes that follows is inside the counted span.\par\medskip
\ifdefined\hypertarget\hypertarget{sourcepointer-2}{}\fi
\noindent\textit{Source reference.} Original lemma 00PG: sources/stacks-chow/tags.txt, source lines 930--930. The reference gives Noetherianity of the integral closure of a one dimensional Noetherian domain in a finite extension, applied at proof lines 35948-35952 to each quotient; the identification of the product with the integral closure is inside the counted span.\par\medskip
\ifdefined\hypertarget\hypertarget{sourcepointer-3}{}\fi
\noindent\textit{Source reference.} Original lemma 031F: sources/stacks-chow/tags.txt, source lines 3982--3982. The reference supplies the finiteness of the residue field extensions for one dimensional rings, applied at proof lines 35952-35955; it is the last input and the conclusion drawn from it is inside the counted span.\par\medskip
\begin{lemma}
\label{lemma-big-extension-is-ok}
Let $A \to B$ be an extension of discrete valuation rings with fraction
fields $K \subset L$. Assume $B$ is essentially of finite type over $A$.
Let $K'/K$ be an algebraic extension of fields such that
the integral closure $A'$ of $A$ in $K'$ is Noetherian. Then the integral
closure $B'$ of $B$ in $L' = (L \otimes_K K')_{red}$ is Noetherian
as well. Moreover, the map $\Spec(B') \to \Spec(A')$
is surjective and the corresponding residue field extensions are finitely
generated field extensions.
\end{lemma}

\begin{proof}
Let $A \to C$ be a finite type ring map such that $B$ is a localization of
$C$ at a prime $\mathfrak p$. Then $C' = C \otimes_A A'$ is a finite type
$A'$-algebra, in particular Noetherian. Since $A \to A'$ is integral, so
is $C \to C'$. Thus $B = C_\mathfrak p \subset C'_\mathfrak p$ is
integral too. It follows that the dimension of $C'_\mathfrak p$ is $1$
(Algebra, Lemma \ref{algebra-lemma-integral-sub-dim-equal}).
Of course $C'_\mathfrak p$ is Noetherian.
Let $\mathfrak q_1, \ldots, \mathfrak q_n$ be the minimal primes
of $C'_\mathfrak p$. Let $B'_i$ be the integral closure of
$B = C_\mathfrak p$, or equivalently by the above of $C'_\mathfrak p$
in the field of fractions of $C'_{\mathfrak p'}/\mathfrak q_i$.
It follows from Krull-Akizuki
(Algebra, Lemma \ref{algebra-lemma-krull-akizuki} applied
to the finitely many localizations of $C'_\mathfrak p$ at its
maximal ideals) that each $B'_i$ is Noetherian.
Moreover the residue field extensions in $C'_\mathfrak p \to B'_i$
are finite by
Algebra, Lemma \ref{algebra-lemma-finite-extension-residue-fields-dimension-1}.
Finally, we observe that $B' = \prod B'_i$ is the integral closure of $B$ in
$L' = (L \otimes_K K')_{red}$.
\end{proof}

\end{document}
